\chapter{Conclusiones}\label{chap:conclusiones}

El objetivo general del Proyecto es una plataforma de análisis de portafolios multiactivo, con clientes web y móvil sobre una misma cuenta, que relacione la composición de la cartera con contexto geopolítico y señales emergentes y presente observaciones cuyo fundamento pueda inspeccionarse. Allocia cumple ese objetivo en su núcleo: una cartera cargada desde cualquiera de los dos clientes se clasifica, se mide y produce observaciones que muestran la regla, los datos y la fuente que las originan. Todas las señales del diseño, incluido el índice GPR, se consultan en vivo. El cumplimiento es parcial en la comprensión de las observaciones por personas sin experiencia: en las sesiones con cinco participantes, dos explicaron una observación sin ayuda, y el efecto de los cambios derivados, ya desplegados, no se volvió a medir. La Tabla~\ref{tab:objetivos-resultado} contrasta el objetivo general y cada objetivo específico de la Subsección~\ref{sec:objetivo} con el resultado y su evidencia.

\begin{table}[ht]
\caption{Cumplimiento del objetivo general y de los objetivos específicos. Fuente: elaboración propia.}\label{tab:objetivos-resultado}
\centering\footnotesize\renewcommand{\arraystretch}{1.2}
\begin{tabular}{>{\raggedright\arraybackslash}p{4.2cm}>{\raggedright\arraybackslash}p{2.8cm}>{\raggedright\arraybackslash}p{7.2cm}}\hline
\textbf{Objetivo} & \textbf{Resultado} & \textbf{Evidencia y límite}\\\hline
\textbf{General:} plataforma web y móvil que vincule la cartera con contexto y señales y explique sus observaciones & Cumplido en el núcleo; parcial en comprensión por perfiles sin experiencia & Flujo completo de cuenta, cartera, métricas y observaciones en ambos clientes (Subsección~\ref{sec:validacion-publicada} y Subsubsección~\ref{sec:reverificacion-octubre13}); seis señales en vivo, incluido el GPR; dos de cinco participantes explican una observación sin ayuda\\
Calcular pesos, exposiciones y concentración con diferencia máxima de $10^{-6}$ & Cumplido en carteras seleccionadas & CP01 a CP05 e IT07; reconstrucción manual de CP03 (Tabla~\ref{tab:lectura-cp03}); no valida la calidad económica de los precios\\
Incorporar contexto geopolítico y señales emergentes con origen, fecha y condición & Cumplido & Índice GPR, Fear \& Greed, dominancia de BTC, TVL, mercados de predicción y noticias con procedencia por componente en el servicio publicado (Tabla~\ref{tab:estado-funcional})\\
Generar observaciones con regla, métricas y evidencia en el 100\,\% de los casos & Cumplido en los casos ejecutados & CP16, IT08 y UI04; la comprensión se trata en el objetivo de evaluación con usuarios\\
Ofrecer el flujo completo desde web y móvil sobre la misma cuenta, con una matriz de paridad & Cumplido; paridad parcial & Prueba de cuenta compartida (Subsubsección~\ref{sec:persistencia-web-movil}); quedan las notificaciones del navegador y algunas vistas y exportaciones de la web\\
Evaluar con casos trazables y con tareas de usuario & Cumplido; criterio de comprensión no alcanzado & Casos trazables en el Anexo~\ref{annex:casos-prueba}; cinco sesiones del Anexo~\ref{annex:protocolo-usuarios}, dos de cinco sin ayuda en la tarea de explicación, ocho cambios derivados\\
Proteger los datos personales y delimitar el uso informativo & Cumplido en lo implementado y probado & Aislamiento entre cuentas, cifrado y baja de cuenta (Tabla~\ref{tab:controles}); sin definir el responsable, los plazos de conservación, la evaluación de la transferencia internacional ni la inscripción de la base\\\hline
\end{tabular}
\end{table}

\section*{Aportes}

El aporte principal es un recorrido inspeccionable desde las posiciones hasta la observación: pesos, exposiciones y concentración se calculan con reglas declaradas, y cada observación muestra su condición, sus métricas y su fuente. Esa decisión responde a dos hallazgos de la investigación. La encuesta ubica la explicación de las observaciones entre las funciones más valoradas, y los reguladores de valores señalan la opacidad de los modelos de lenguaje y el sesgo de automatización como riesgos del asesoramiento automatizado. Allocia separa, por eso, el motor determinista que explica de la capa de lenguaje que redacta sobre fuentes citadas, y mantiene la decisión en la persona.

El segundo aporte es la integración, en un mismo flujo y en español, de activos tradicionales y criptoactivos con contexto externo fechado. Las herramientas relevadas resuelven partes de esa tarea; ninguna documenta el conjunto. El tercero es la evidencia reproducible: 62 comprobaciones en la evaluación inicial, repetidas en la reverificación del 13 de octubre junto con 33 pruebas unitarias, 27 de la suite Playwright y 14 comprobaciones web adicionales, que detectaron y corrigieron un defecto en la baja de cuenta y dos de accesibilidad.

\section*{Limitaciones}

Los resultados valen para los casos ejecutados y el entorno documentado. La corrección aritmética es independiente de la calidad económica de las cotizaciones, y el HHI es una medida de riesgo incompleta porque omite correlaciones y volatilidad. La cobertura de datos depende de proveedores con cuotas: las 34 cotizaciones del catálogo y los análisis reales con Anthropic se obtuvieron en momentos registrados, no de forma continua. La investigación con usuarios se apoya en una muestra de conveniencia, tres entrevistas transcriptas y tres personas con atributos hipotéticos. Las sesiones con cinco participantes muestran que la explicación resulta clara para inversores activos e insuficiente para quien recién empieza; son cinco experiencias por conveniencia, sin valor estadístico. El sistema no está auditado por terceros, no está inscripto como agente ante la CNV y no debe presentarse como asesoramiento.

\section*{Trabajo futuro}

Los ocho cambios derivados de las sesiones (rótulos de fechas, distinción entre observación de cartera y contexto, explicación simple de la concentración, edición y baja de cuenta más claras, confirmaciones persistentes) están implementados y desplegados; la primera línea de trabajo es repetir la evaluación con perfiles ocasionales y principiantes. En el producto siguen las notificaciones en el navegador y el cierre de la matriz de paridad. Como evolución, la descomposición de fondos en sus tenencias, medidas de riesgo adicionales a la concentración y la validación financiera de las reglas con series históricas, siempre dentro de la finalidad informativa del Proyecto. La viabilidad del modelo de negocio depende de verificar la disposición a pagar relevada en la encuesta con usuarios reales, condición que ninguna prueba técnica reemplaza.
